\section{带标记的完备化与反射行能量}
\label{sec:marked}\label{sec:marked-energy}

补偿的低估计与逆矩使用同一个带标记的完备化。我们在此建立其反射均方。低估计将使用没有额外无平方因子标记或穿刺的情形；更一般的形式允许由逆矩的 Poisson 变换所引入的标记与掩码。

我们沿用节~\ref{sec:arithmetic} 的记号。特别地，理想变量取其乘性主生成元，固定排除集 \(S\) 包含 \(6\) 之上的素理想，且剩余符号的任意幂在非单位元上为零，即使其指数可被六整除。元素行不必无平方因子，也不必与 \(S\) 互素。

\subsection{完备化和与积形式的标记}

固定一个有限指标集 \(I\) 以及 \(S\) 之外素理想的配对不交列表 \(\mathcal P_i\)，其中对 \(p\in\mathcal P_i\) 有 \(q_p\asymp Z^{z_i}\)。这些列表及其有界的单个系数 \(a_i(p)\) 独立于当前行与无平方因子标记。它们的名义长度 \(z_i\) 落在固定有界的区间内。

此处及下文，\emph{固定穿刺} 指一个函数
\[
 \rho(n)=1_{(n,\mathfrak r_\rho)=1},
\]
其中 \(\mathfrak r_\rho\) 无平方因子，且仅要求在所指示的当前行与标记求和内固定。此类函数的有界乘积具有同样形式，其中 \(\mathfrak r_\rho\) 等于其模之积的根式。该根式的范数将被显式地界住。它可以依赖于先前固定的外部理想，但不依赖于当前的任一平均变量。

设 \(f\) 为 \(S\) 之外的无平方因子理想，其中对固定有界区间内的非负 \(V\) 有 \(q_f\asymp Z^V\)，并设 \(k\in\mathcal O\)。对一个固定的乘性有限射线特征 \(\nu\)，其完整的零延拓定义模的素支集在 \(S\) 中，仅对 \(S\) 之外的无平方因子主理想 \(n\) 定义 \(a(n)\)，并对 \(S\) 之外的每个主理想 \(n\) 通过下式定义 \(\Psi_k(n)\)
\[
 a(n)=\bar\alpha(n)\gamma_2(n)\nu(n)\rho(n),\qquad
 \Psi_k(n)=\nu(n)\rho(n)\chi_n(k)\chi_n(f)^4.
\]
有限特征、穿刺以及所有槽列表均独立于 \(k,f\)。有限射线群上的固定函数在使用此定义之前总是展开为真正的群特征；因此，即使先前某一步产生了特征的有限线性组合，\(\nu\) 仍是乘性的。

对具有有界系数的配对不交素理想列表的子集 \(I'\subset I\)，通过下式定义其标记
\begin{equation}\label{old-eq:4.2}
 \mathfrak d_{I'}(A)=
 \sum_{(p_i)\in\prod_{i\in I'}\mathcal P_i}
      \prod_{i\in I'}a_i(p_i)1_{p_i\mid A}.
\end{equation}
当 \(I'\) 不言自明时，我们将其从记号中省略。由于这些列表配对不交，一个元组具有无平方因子的乘积。对每个固定的 \(\epsilon>0\)，整除 \(A\) 的元组数目为 \(O_\epsilon(q_A^\epsilon)\)。方程~\eqref{old-eq:4.2} 中的系数是个别素理想函数之积；当某些槽被分配到某个提取出的因子时，所利用的正是这一性质，而不仅仅是逐点的除子界。

对一个光滑环形 \(W\)，置 \(V_*(y)=y^{1/2}W(y)\)，并定义完备化带标记和
\begin{equation}\label{eq:marked-completion}
 \mathcal T_{\mathfrak d}(X;k)=
 \sum_{\substack{n\ {\rm sf}\\b}}
 \frac{\bar\alpha(n)\gamma_2(n)\Psi_k(n)\,
       \bar\alpha(b)^3\Psi_k(b)^3}{\sqrt{q_n}\,q_b}
 \mathfrak d(nb^3)V_*(q_nq_b^3/X).
\end{equation}
两个原始理想都避开 \(S\)；它们可以共享素理想。标记属于整个指标 \(nb^3\)。当 \(\mathfrak d=1\) 时，这就是方程~\eqref{eq:reflection} 左侧的 Mellin 完备化。当只使用完备化估计时，来自该方程的 Gaussian Mellin 检验也是允许的。

对于反射公式，我们使用引理~\ref{lem:reflection-row-sectors}，仅对一个非零行。如该引理中那样写 \(k=u k_S k_{\rm good}\)。在固定单位、模六的 \(S\)-赋值以及好行射线类之后，其字面扇形特征为
\[
 \Psi_0^{[k]}(A)=
 \begin{cases}
 \nu(A)\chi_A(u k_S)\mathcal R(A,k_{\rm good}),
       &A\equiv1\pmod3,\ (A,S)=1,\\
 0,&\text{otherwise}.
 \end{cases}
\]
零分支在任一特征之前求值。该引理在好素理想变动之前提供了一个有限族，从而提供了一个共同的 \(L\)，其素支集为 \(S\)。写 \(M_{\rm ref}=\lambda^{12}L^4\)；完整的反射扇形模为 \(M_{\rm ref}^2\)。\(k,f,\mathfrak r_\rho\) 的所有好素理想都保留在局部素理想集中，其指数按模六合并，且所有零都保留。\(f\) 处的四次方不产生额外的互反符号，因为 \(\mathcal R(A,f)^4=1\)。穿刺是其好素理想处局部零幂之积。因此，局部的表述在所声明的原始支集上恰好给出 \(\Psi_k(n)\Psi_k(b)^3\)，包括在共享的好素理想处；没有变动的好素理想进入 \(L\)。

\subsection{整指标的带标记反射}\label{par:reflection-marks}

下一个推论把积形式的标记加入到第 I 部分中的精确反射之中。其分支数据很重要：一个非活跃的带标记素理想不出现在导子中，因此带标记的变换不能通过将一个固定导子的对偶和乘以独立的局部因子而得到。

\begin{corollary}[带整指标标记的完备化反射]
\label{cor:marked-reflection}
固定命题~\ref{prop:completed-reflection} 的一次调用，其基局部素理想集为 \(\mathcal P_0\)、幂为 \(j_p\)、固定乘子为 \(\phi\)、检验为 \(V\)。对此调用置 \(M_{\rm ref}=\lambda^{12}L^4\)。设 \(\mathcal Q\) 为与 \(\mathcal P_0\) 不交的互异好素理想的有限集。在直接的完备化系数和中，插入因子 \(\prod_{q\in\mathcal Q}1_{q\mid nb^3}\)；将此和记作 \(\mathcal T_{\mathcal Q}(X;\Psi,V)\)。它与方程~\eqref{eq:marked-completion} 具有相同的正规化，其中 \(\Psi_k,V_*\) 被替换为 \(\Psi,V\)。原始的无平方因子理想与立方理想可以共享素理想。

存在一个精确的有限展开
\begin{equation}\label{eq:marked-reflection}
 \begin{split}
 \mathcal T_{\mathcal Q}(X;\Psi,V)
 =\sum_{\mathcal B}\zeta_{\mathcal Q}
 \sum_{0\ne\mu\in\lambda^{-4}\mathcal O}
 &\frac{d(\mu)\alpha(\mu)\vartheta(\lambda^4\mu)}{\sqrt{q_\mu}}
 \prod_{p\in\mathcal A_0}B_p(\lambda^4\mu)\\
 &\cdot\prod_{q\in\mathcal A_{\mathcal Q}}
       q_q^{-1/2}\chi_q(\lambda^4\mu)^{-2}
 V^\sharp\!\left(\frac{q_\mu X}{q_c^2}\right).
 \end{split}
\end{equation}
这里一个分支 \(\mathcal B\) 指定 \(h_0\bmod L\)、基活跃集与非活跃集 \(\mathcal A_0,\mathcal I_0\)，以及一个划分 \(\mathcal Q=\mathcal A_{\mathcal Q}\sqcup\mathcal I_{\mathcal Q}\)。所有满足 \(j_p\ne0\) 的基素理想都是活跃的。置
\[
 r=\prod_{p\in\mathcal A_0\cup\mathcal A_{\mathcal Q}}p,
 \qquad c=c_F(h_0)r.
\]
基因子 \(B_p\) 以及典范函数 \(d,\vartheta\) 就是命题~\ref{prop:completed-reflection} 对此活跃根式所给出的那些。当每个 \(\sigma_p,\epsilon_p,\omega_{p,j}\) 都使用该分支的整个 \(c\) 计算时，标量为
\begin{equation}\label{eq:marked-reflection-phase}
 \begin{split}
 \zeta_{\mathcal Q}={}&-\frac{i}{81}\bar\alpha(c)^2
       \widehat\phi(h_0)\bar\kappa_F
       \prod_{p\in\mathcal I_0}(1-q_p^{-1})
       \prod_{q\in\mathcal I_{\mathcal Q}}q_q^{-1}\\
 &\cdot\prod_{p\in\mathcal A_0}\chi_p(\sigma_p)^{-2}\omega_{p,j_p}
       \prod_{q\in\mathcal A_{\mathcal Q}}
             \chi_q(\sigma_q)^{-2}(-\omega_{q,0}).
 \end{split}
\end{equation}
特别地 \(|\zeta_{\mathcal Q}|\le1/81\)。函数 \(d,\vartheta\) 与标量 \(\kappa_F\) 具有源自第 I 部分的、对 \(h_0\) 与 \(r\bmod M_{\rm ref}^2\) 的典范依赖性。每个对偶和都绝对收敛，并且与那里具有相同的完整源支集。若某个带标记素理想反而属于基局部集，则相应的原始项恒为零，并在使用此公式之前被移除。
\end{corollary}

\begin{proof}
对每个好素理想和每个 \(A\in\mathcal O\)，零约定给出逐点恒等式 \(1_{q\mid A}=1-\chi_q(A)^0\)。将它在 \(\mathcal Q\) 上的乘积展开，并对所得到的有限多个完备化和中的每一个应用命题~\ref{prop:completed-reflection}。对固定的 \(h_0\) 与活跃根式，无论一个缺失的素理想是从局部集中被省略，还是以其非活跃分支中指数零的形式被纳入，该命题都使用相同的 \(c_F,c,d,\vartheta,\kappa_F\)、其他局部相位以及核。于是 \(q\) 处的两个非活跃贡献以系数 \(1-(1-q_q^{-1})=q_q^{-1}\) 合并。活跃贡献仅来自被减去的零掩码变换，其局部因子为 \(q_q^{-1/2}\chi_q(\lambda^4\mu)^{-2}\)，标量为 \(-\omega_{q,0}=\tau_{q,2}^{+}\chi_q(\epsilon_q)^{-2}\)。在所有带标记素理想处汇集这些选择即证明了方程~\eqref{eq:marked-reflection} 与 \eqref{eq:marked-reflection-phase}。非活跃素理想不出现在 \(r\) 中，从而不出现在该分支的导子与核尺度中。这是相容分支之间的有限容斥，而不是一个新的 theta 恒等式。

若某个带标记素理想是基局部素理想，则该标记迫使它整除 \(nb^3\)，而此时该素理想处的零延拓基因子为零，包括指数为零的情形。这就证明了最后一个断言，而无需对零符号作任何除法。它还涵盖了原始的无平方因子理想与立方理想共享该带标记素理想的情形。
\end{proof}

对于行范数，本推论的相位推论必须在其精确范围内陈述。对互异的活跃好素理想 \(p,q\)，方程~\eqref{eq:reflection-cross-prime-phase} 与三次互反律给出
\begin{equation}\label{eq:reflection-pair-phase}
 \chi_p(q)^{2j_p+2}\chi_q(p)^{2j_q+2}
       =\left(\frac qp\right)_3^{j_p+j_q+2}.
\end{equation}
对一个活跃标记，指数为 \(j_q=0\)；其符号的改变仅出现在其单素理想标量中。因此，一个 \(j_p=1\) 的剩余行素理想与一个带标记素理想的配对因子为 \((q/p)_3^3=1\)。而一个 \(j_p=1\) 的行素理想与一个变动的 \(j_q=4\) 素理想则会留下 \((q/p)_3\)，它未必等于一。因此所有非剩余基素理想，尤其是每个 \(j=4\) 素理想，都在内层行范数之前固定。

记 \(R\) 为剩余 \(j=1\) 行素理想之积，\(P\) 为活跃带标记素理想之积。原始的行-标记冲突首先由推论的最后一个断言移除，随后由 \(1_{(P,R)=1}\) 表示。若冻结的活跃基乘积为 \(F_{\rm act}\)，则 \(r=F_{\rm act}RP\)。分别固定 \(R\) 与 \(P\) 模完整 \(M_{\rm ref}^2\) 的类，而不仅仅是它们乘积的类，也不仅仅是它们较小的互反类。这样就在不添加配对相关限制的情况下固定了尖点数据。行-行相位是行标量，标记-标记相位是元组标量，它们与冻结素理想的相互作用仅依赖于这些集合中的一个。角向导子标量以同样的方式分解。

在固定对偶单位、其 \(\lambda\)-赋值以及提取出的冻结因子之后，将剩余的对偶部分写作 \(nb^3\)。变动的列恰好为
\[
 \chi_R(nb^3)^3=\chi_R(nb)^3,\qquad
 \chi_P(nb^3)^{-2}=\chi_P(n)^{-2}1_{(P,b)=1}.
\]
这些恒等式包含每一个零。由某个行或某个元组与某个提取出的冻结因子相遇所造成的零，是对该行或元组的固定限制。源系数、固定加性因子以及所有冻结的局部列因子都是完整提取对偶指标的函数。因此，在共同的光滑分离之后，元组系数可以是元组的、独立于 \(R,n,b\) 的一般有界函数，而对偶系数独立于 \(R,P\)。当槽被分配到提取出的因子时，标记先前的积形式仍然是必需的。

\subsection{一个二次--三次范数估计}

第 I 部分确立了 Goldmakher--Louvel 的二次大筛法所需的方向，并在引理~\ref{lem:completed-quadratic-reduction} 中证明了保零的完备化归约。我们加入 Heath--Brown 的三次大筛法：
\begin{equation}\label{eq:cubic-large-sieve}
 \sum_{\substack{a\equiv1\ (3)\\q_a\le U}}^{*}
 \left|\sum_{\substack{b\equiv1\ (3)\\q_b\le V}}^{*}
           c_b\left(\frac ba\right)_3\right|^2
 \ll_\epsilon (UV)^\epsilon
       \{U+V+(UV)^{2/3}\}\sum_b^{*}|c_b|^2,
\end{equation}
其中两个星号都表示在 \(\mathcal O\) 中无平方因子，且系数是任意复数 \cite[定理 2]{HB}。指标不必具有无平方因子的有理范数，也不排除有理素因子。共轭与三次互反律允许相反的方向。在此筛法与方程~\eqref{eq:quadratic-large-sieve} 中，对行集的固定限制使正的外部和减小。而对系数支集的固定限制则通过将省略的系数置零并应用同一定理来实现；不声称列支集的单调性。这一观察不允许任意的配对相关掩码；下一个引理解它所使用的两个此类掩码。

\begin{lemma}[一个二次--三次范数界]\label{lem:hybrid-energy}
设 \(K,N,B,L\ge1\)。设 \(k,n,P\) 为无平方因子主理想，满足 \(q_k\ll K\)、\(q_n\asymp N\)、\(q_P\asymp L\)；设 \(b\) 为任意主理想，满足 \(q_b\asymp B\)。理想 \(k,P\) 避开当前的排除集 \(S\)；允许额外的固定排除。源理想 \(n,b\) 可以含有 \(S\) 中除 \(\lambda\) 之外的素理想。设 \(a(P)\) 与 \(\beta(n,b)\) 满足
\[
 |a(P)|\le1,\qquad |\beta(n,b)|\le1.
\]
这里 \(a\) 独立于 \(k,n,b\)，\(\beta\) 独立于 \(k,P\)。则
\begin{equation}\label{old-eq:4.6b}
 \begin{split}
 &\sum_k\left|
 \sum_P\frac{a(P)}{\sqrt{q_P}}1_{(P,k)=1}
 \sum_{n,b}\beta(n,b)\chi_k(nb)^3
                 \chi_P(n)^{-2}1_{(P,b)=1}
 \right|^2\\
 &\qquad\ll_\epsilon
 (KNBL)^\epsilon(K+NB)B\{N+L+(NL)^{2/3}\}.
 \end{split}
\end{equation}
允许对 \(k\)-集的固定限制，以及对 \((n,b)\)-支集或 \(P\)-支集的固定限制，也允许固定的射线扇形。所有此类限制都必须保持上述两个系数独立性条件。若除子有界的重数保持这些条件，则可以包含它们；否则必须先用证明中的除子 Cauchy 不等式将其移除。
\end{lemma}

\begin{proof}
将方程~\eqref{old-eq:4.6b} 的左侧记作 \(\mathcal H\)。首先移除完整的变动行-标记掩码：
\[
 1_{(P,k)=1}=\sum_{D\mid(P,k)}\mu(D).
\]
对每一行，此类除子的数目至多为 \(d_{\mathcal O}(k)\)。逐行的除子 Cauchy 不等式，随后是对行取正的和，因此花费一个允许的因子 \((KNBL)^\epsilon\)。在固定 \(D\)-被加项中写 \(P=DP'\) 与 \(k=Dk'\)。无平方因子性给出固定限制 \((D,P')=(D,k')=1\)。定义
\[
 \begin{split}
 a_D(P')&=a(DP')1_{(D,P')=1},\\
 \beta_D(n,b)&=\beta(n,b)\chi_D(nb)^3\chi_D(n)^{-2}1_{(D,b)=1}.
 \end{split}
\]
所有这些因子都保留它们的零，且 \(|\beta_D|\le1\)。将 \(q_P^{-1/2}=q_D^{-1/2}q_{P'}^{-1/2}\) 分解出来给出
\begin{equation}\label{eq:hybrid-mark-mask-reduction}
 \begin{split}
 \mathcal H\ll (KNBL)^\epsilon\sum_D\frac1{q_D}
 \sum_{k'}\left|
 \sum_{n,b}\beta_D(n,b)\chi_{k'}(nb)^3
 \sum_{P'}\frac{a_D(P')}{\sqrt{q_{P'}}}
       \chi_{P'}(n)^{-2}1_{(P',b)=1}
 \right|^2.
 \end{split}
\end{equation}
这里 \(q_{k'}\ll K/q_D\)、\(q_{P'}\asymp L/q_D\)，且所有继承的固定限制都被保留。外层的 \(D\) 是该分解所允许的无平方因子好除子，因此 \(q_D\ll\min(K,L)\)。特别地，在单个 \(D\)-被加项内部不会恢复任何 \((P',k')\) 掩码。其抵消属于已经被逐行 Cauchy 界定过的完整 M\"obius 和。

对固定 \(D\)，\(\beta_D(n,b)\) 与内层 \(P'\)-和的乘积是一个独立于 \(k'\) 的任意复系数。以行长 \(K/q_D\) 应用方程~\eqref{eq:completed-quadratic-reduction}。使用其记号
\[
 b=gt^2,\qquad c=(n,g),\qquad n=cm,\qquad g=ch,
 \qquad q_c\asymp C,\quad q_g\asymp G,\quad q_t\asymp T,
\]
于是 \(B\asymp GT^2\)，并对其完整的 \(c\)-掩码除子写 \(c=rc'\)。理想 \(c,m,h\) 两两互素且无平方因子；不施加 \(g\) 与 \(t\) 之间的互素性。所引用的归约已经包含了自然的 \(k'\)-\(t\) 限制以及 \(k'\)-\(c\) 掩码的完整 M\"obius 展开。它不会在剩余的行与 \(c'\) 之间恢复掩码。

对一个 \(C,G,T\) 块，其插入到方程~\eqref{eq:hybrid-mark-mask-reduction} 中的正输出至多为
\begin{equation}\label{eq:hybrid-after-quadratic-reduction}
 \begin{split}
 (KNBL)^\epsilon\sum_D\frac1{q_D}T\sum_{q_t\asymp T}
 \sum_{\substack{r\ {\rm sf},\ (r,S)=1\\q_r\ll\min(K/q_D,C)}}
 &\left\{\frac K{q_Dq_r}+\frac{NG}{C^2}\right\}\frac C{q_r}
 \sum_{c',m,h}|\beta_D(rc'm,rc'ht^2)|^2\\
 &\quad\cdot\left|
 \sum_{P'}\frac{a_{D,r}(P')}{\sqrt{q_{P'}}}
       \chi_{P'}(c'm)^{-2}1_{(P',ht)=1}
 \right|^2,
 \end{split}
\end{equation}
其中正的 \(c',m,h\)-和保持原支集，并且
\[
 a_{D,r}(P')=a_D(P')\chi_{P'}(r)^{-2}.
\]
这里所使用的分解，包含所有零，为
\[
 \chi_{P'}(rc'm)^{-2}1_{(P',rc'ht^2)=1}
 =\chi_{P'}(r)^{-2}\chi_{P'}(c'm)^{-2}1_{(P',ht)=1}.
\]
特征已经提供了 \(r,c'\) 与 \(m\) 处缺失的零。因此 \(r\)-因子只与元组有关，包括当 \((P',r)\ne1\) 时它的零。固定的 \(D\)-相位在此正界之前一直留在 \(\beta_D\) 内。二次归约在对 \(m,h\) 的好乘积进行筛分之前，已经分离出它们有限多个无平方因子的 \(S\)-部分；方程~\eqref{eq:hybrid-after-quadratic-reduction} 中的 \(m,h\) 是原来的理想。

我们现在可以舍弃 \(|\beta_D|^2\le1\)。对固定的 \(h,t,D,r\)，\(P'\)-系数独立于 \(j=c'm\)，理想计数给出
\[
 \sum_{P'}\frac{|a_{D,r}(P')|^2\,1_{(P',ht)=1}}{q_{P'}}
 \ll_\epsilon (KNBL)^\epsilon.
\]
乘积 \(j=c'm=n/r\) 是无平方因子主理想，其范数为 \(O(N/q_r)\)。它可以含有 \(S\) 中除 \(\lambda\) 之外的允许素理想，因为只有二次的列被剥去其固定的 \(S\)-部分。按 \(j\) 对 \(c',m\) 分组至多花费一个除子因子。只扩大这个正的无平方因子 \(j\)-区间。三次互反律、共轭以及方程~\eqref{eq:cubic-large-sieve}，随后是 \(O(G/C)\) 个 \(h\) 的选择，在 \(q_D\asymp D_0,\ q_r\asymp r_0\) 上给出
\begin{equation}\label{old-eq:4.6c}
 (KNBL)^\epsilon\frac GC
 \left\{\frac N{r_0}+\frac L{D_0}
       +\left(\frac{NL}{r_0D_0}\right)^{2/3}\right\}.
\end{equation}
\(ht\) 处的固定掩码在此三次筛法中是一个系数限制，而不是行相关的掩码。

该归约提供 \(T\sum_t\)，其中 \(t\) 有 \(O(T)\) 种可能，并提供因子 \(C/q_r\)。与 \(h\)-计数一起，它们贡献
\[
 T^2(C/r_0)(G/C)\asymp B/r_0.
\]
在其二进区间中有 \(O(r_0)\) 个可能的 \(r\)，而在其区间上 \(q_D^{-1}\) 之和为 \(O(1)\)，至多相差一个允许的小幂。当 \(D_0,r_0\) 被替换为 \(1\) 时，两个筛因子都增大，且 \(NG/C^2\ll NB\)。对对数多个二进区间求和即证明了方程~\eqref{old-eq:4.6b}。亚单位商区间为空，有界区间通过改变固定的环形常数被纳入。此论证既未引入 \((g,t)=1\) 条件，也未引入新的配对相关掩码。
\end{proof}

\subsection{反射能量界}

我们接下来量化方程~\eqref{eq:marked-completion} 的反射。下面的描述也规定了陈述中所使用的二进贡献。

\paragraph{反射块数据.}\label{par:reflected-energy-data}
固定 \(f,\rho\) 并置 \(Q=\log_Zq_{\mathfrak r_\rho}\)。对一个元素行 \(0<q_k\ll Z^M\)，通过下式定义其理想的极大幂强部分
\[
 k_{\rm pow}=\prod_{v_p(k)\ge2}p^{v_p(k)}.
\]
将赋值一的素理想分成支集在 \(S\cup\operatorname{supp}(f\mathfrak r_\rho)\) 上的无平方因子乘积 \(k_{\rm sup}\)，以及该集合之外的无平方因子乘积 \(k_{\rm res}\)。这三个素理想支集两两不交，且它们的乘积是 \(k\) 的理想。此处一个幂强理想意指每个正的素理想赋值至少为二。固定 \(k\) 的单位。在
\[
 q_{k_{\rm pow}}\asymp Z^O,\qquad q_{k_{\rm res}}\asymp Z^H.
\]
中插入光滑二进划分。选择这些理想中心为非负，单位二进区间的中心为零。特别地，一个实际的环形因子 \(\chi_{\rm row}(q_{k_{\rm res}}/Z^H)\) 被保留在行和中。有界二进区间包含 \(k_{\rm res}=1\)。从这些支集中选取固定常数 \(C_k,C_O,C_H\ge1\)，使得
\[
 q_k\le C_kZ^M,\qquad q_{k_{\rm pow}}\ge Z^O/C_O,\qquad
 q_{k_{\rm res}}\ge Z^H/C_H.
\]
由于 \(q_{k_{\rm sup}}\ge1\)，一般范数不等式方程~\eqref{eq:actual-residual-row-dyad} 在此给出
\[
 H\le M-O+\frac{\log C_{\rm res}}{\log Z},
 \qquad C_{\rm res}=C_HC_kC_O.
\]
此不等式不必是等号，因为 \(k_{\rm sup}\) 可能具有正的范数长度。

冻结实际的理想 \(k_{\rm pow},k_{\rm sup}\)。在 \(S\) 之外的每个素理想处，将 \(\rho(n)\chi_n(k)\chi_n(f)^4\) 的局部指数按模六合并，始终保留零延拓。\(k_{\rm res}\) 的素理想具有指数 \(j=1\)。所有其他此类素理想，包括每个指数 \(j=4\) 的非槽素理想，现在都被固定。在反射公式中，当一个素理想出现在 \(c\) 中时称其为活跃的。对活跃零指数素理想的一个固定选择以及每个 \(j=4\) Ramanujan 因子的一个固定分裂，令：
\begin{center}
\begin{tabularx}{\linewidth}{@{}>{$}l<{$} X@{}}
 A_0 & 为 \(k_{\rm res}\) 之外所有活跃非槽素理想的对数范数,\\
 S_0 & 为其中小 Ramanujan 项与活跃零掩码项的对数范数,\\
 N_0 & 为分配到无平方因子对偶理想的 \(j=4\) 整除项的对数范数,\\
 B_0 & 为分配到立方理想但不分配到无平方因子理想的 \(j=4\) 整除项的对数范数.
\end{tabularx}
\end{center}
这里一个对数范数是指所示乘积范数的 \(\log_Z\)。该分配使用
\(1_{p\mid nb^3}=1_{p\mid n}+1_{p\nmid n}1_{p\mid b}\)。
设 \(z_a=\sum_{i\ {\rm active}}z_i\) 为其名义槽长度之和。活跃素理想之积的范数除以 \(Z^{z_a}\) 落在一个固定的紧区间内，包括名义长度为零或空列表的情形。在提取出被迫的无平方因子素理想与立方素理想之后，将剩余的对偶理想限制为
\[
 q_n\asymp Z^v,\qquad q_b\asymp Z^{\ell_b},\qquad
 q_{\lambda^{h_\lambda}}=Z^{e_\lambda}.
\]
选择理想中心 \(v,\ell_b\) 为非负。这些是源对偶理想：它们保留方程~\eqref{old-eq:4.3} 所允许的每个 \(S\) 素理想；不在其上放置额外的 \(S\)-掩码。只施加所显示的 \(\lambda\)-赋值支集。单位以及整数 \(h_\lambda\ge-4\) 在此二进区间中被固定。对这些选择，记 \(\mathcal U_{v,\ell_b,e_\lambda}(k_{\rm res})\) 为方程~\eqref{eq:marked-reflection} 的右侧，连同积形式的元组系数求和，并带上所指定的局部选择与对偶限制，以及实际因子 \(\chi_{\rm row}(q_{k_{\rm res}}/Z^H)\)。在此定义中活跃槽元组仍被求和。此定义对每个固定的、幂强且被支撑的行部分，以及对通过固定剩余行乘积与活跃槽乘积模完整反射模 \(M_{\rm ref}^2\) 的各自类而得到的每个扇形，分别作出。

\begin{lemma}[反射能量]\label{lem:reflected-energy}
沿用上文的 \hyperref[par:reflected-energy-data]{反射块数据}。设 \(Z\ge2\)，并假设它们的对数长度在固定有界集合中变动。在方程~\eqref{eq:marked-completion} 的完备化和中，固定的特征、穿刺与槽系数独立于 \(k,f\)，槽支集不交，且槽系数如方程~\eqref{old-eq:4.2} 那样为乘积。在冻结 \(k_{\rm pow},k_{\rm sup}\) 之后的固定行限制必须独立于活跃槽与对偶变量，除了显式的零掩码 \(1_{(P,k_{\rm res})=1}\) 之外。固定一个小的 \(\tau_{\rm ref}>0\)，置 \(X=Z^{N_*}\)。定义 \(T_d\)，并按下式称一个二进区间被保留：
\begin{equation}\label{old-eq:4.4}
 \begin{gathered}
 T_d=2H+2A_0+2z_a-N_*-N_0-3B_0,\\
 v+3\ell_b+e_\lambda\le T_d+\tau_{\rm ref}
 \quad\text{对一个被保留的二进区间}.
 \end{gathered}
\end{equation}
满足
\[
 v+3\ell_b+e_\lambda>T_d+\tau_{\rm ref}
\]
的项具有任意小的总大小，代价是一个固定的多项式高度成本，这要在足够远的核围道移动之后，并且 \(Z\) 的阈值仅依赖于固定的支集区间与 \(\tau_{\rm ref}\)。每个被保留的二进区间，对每个活跃元组，其系数指数为
\begin{equation}\label{old-eq:4.5}
 -v/2-\ell_b-e_\lambda/3-(S_0+B_0+z_a)/2.
\end{equation}
定义
\begin{equation}\label{old-eq:4.7}
 u=\min\{v,z_a,(v+z_a)/3\},
\end{equation}
以及
\begin{equation}\label{old-eq:4.6}
 \begin{split}
 E_{\rm ref}={}&O/2+\max(H,v+\ell_b)-S_0-B_0+z_a-u-\ell_b
                     -2e_\lambda/3\\
 &-\tfrac12(T_d-v-3\ell_b-e_\lambda)_+.
 \end{split}
\end{equation}
对每个 \(\epsilon>0\)，具有固定 \(k_{\rm pow},k_{\rm sup}\) 的贡献满足
\[
 \sum_{k_{\rm res}\ {\rm sf}}
       |\mathcal U_{v,\ell_b,e_\lambda}(k_{\rm res})|^2
 \ll Z^{E_{\rm ref}-O/2+\epsilon}.
\]
对其 \(O\)-二进区间中的固定部分以及局部选择求和，将此变为 \(O(Z^{E_{\rm ref}+\epsilon})\)。这些界在所声明的区间内对变动的理想与穿刺是一致成立的，带有限多个光滑半范数以及范数扭动高度的一个固定多项式成本。当 \(e_\lambda\) 为负时，其值为 \(O(1/\log Z)\)。特别地，该公式仅用于满足方程~\eqref{eq:actual-residual-row-dyad} 的实际剩余行二进区间。
\end{lemma}
